\section*{Capítulo \ref{ch:g12:selection-drift-speciation} --- Selección, deriva y especiación}

\begin{solution}{exo:g12:selection-drift-speciation:1}
La fracción de todas las copias de un \omterm{def:g10:universal-dna:gene}{gen} de una \omterm{def:g12:selection-drift-speciation:population}{población} que son un
\omterm{def:g10:universal-dna:gene}{alelo} dado. Con $p = 0.8$ y $q = 0.2$: $AA$ 64\%, $Aa$ 32\%, $aa$ 4\%.
\end{solution}

\begin{solution}{exo:g12:selection-drift-speciation:2}
Selección: una diferencia de reproducción entre \omterm{def:g11:genes-and-disease:genetic}{portadores} de \omterm{def:g10:universal-dna:gene}{alelos}
distintos, fijada por el ambiente, que cambia las \omterm{def:g12:selection-drift-speciation:population}{frecuencias} en una
dirección. Deriva: fluctuación al azar de las \omterm{def:g12:selection-drift-speciation:population}{frecuencias} por el
muestreo de una \omterm{def:g12:selection-drift-speciation:population}{población} finita, sin ninguna dirección. La selección
depende de lo que hace el \omterm{def:g10:universal-dna:gene}{alelo}; la deriva, no.
\end{solution}

\begin{solution}{exo:g12:selection-drift-speciation:3}
Alrededor del 75\% en 1880, cuando el hollín había ennegrecido los
troncos; alrededor del 30\% en 1980, después de que las leyes de aire
limpio dejaran aclararse otra vez los troncos.
\end{solution}

\begin{solution}{exo:g12:selection-drift-speciation:4}
Los \omterm{def:g10:universal-dna:gene}{alelos} de la generación siguiente son una muestra de los de la
anterior; una muestra pequeña se desvía más del origen que una grande,
así que las \omterm{def:g12:selection-drift-speciation:population}{frecuencias} fluctúan más y un \omterm{def:g10:universal-dna:gene}{alelo} puede perderse por
azar.
\end{solution}

\begin{solution}{exo:g12:selection-drift-speciation:5}
Una interrupción del intercambio de \omterm{def:g10:universal-dna:gene}{genes} entre dos \omterm{def:g12:selection-drift-speciation:population}{poblaciones}, lo
bastante larga para que diverjan hasta que ya no puedan cruzarse.
\end{solution}

\begin{solution}{exo:g12:selection-drift-speciation:6}
Copias de $a$: $160 + 40 = 200$ de 1000, así que $q = 0.2$ y $p = 0.8$.
Esperados: $0.64 \times 500 = 320$, $0.32 \times 500 = 160$, $0.04
\times 500 = 20$ --- exactamente el censo: la \omterm{def:g12:selection-drift-speciation:population}{población} está en el
estado de referencia.
\end{solution}

\begin{solution}{exo:g12:selection-drift-speciation:7}
$q^2 = 0.0001$: uno de cada diez mil muestra el rasgo; $2pq \approx
0.02$: uno de cada cincuenta lo lleva. El noventa y nueve por ciento de
las copias está en heterocigotos, sobre los que la selección contra el
rasgo no tiene ningún agarre.
\end{solution}

\begin{solution}{exo:g12:selection-drift-speciation:8}
Las dos llegaron a un extremo en la generación 12, una fijada en el
100\% y la otra perdida. La tercera vagó sin tocar ninguno de los dos
límites en 20 generaciones --- de nuevo el azar; con más tiempo lo
haría.
\end{solution}

\begin{solution}{exo:g12:selection-drift-speciation:9}
Las mariposas claras siguieron reproduciéndose en el campo sin
contaminar, donde estaban favorecidas, y las mariposas vuelan: los
inmigrantes devolvían el \omterm{def:g10:universal-dna:gene}{alelo} claro a la región con hollín en cada
generación, así que nunca desapareció.
\end{solution}

\begin{solution}{exo:g12:selection-drift-speciation:10}
Veinte individuos llevan como mucho cuarenta copias de cada \omterm{def:g10:universal-dna:gene}{gen}, una
muestra minúscula de los \omterm{def:g10:universal-dna:gene}{alelos} originales; casi todos se perdieron en
el cuello de botella, y la deriva de la \omterm{def:g12:selection-drift-speciation:population}{población} pequeña que vino
después perdió más. Los efectivos volvieron; los \omterm{def:g10:universal-dna:gene}{alelos} no pueden
volver, salvo por \omterm{def:g11:mutations:mutation}{mutación} nueva.
\end{solution}

\begin{solution}{exo:g12:selection-drift-speciation:11}
Diez \omterm{def:g12:selection-drift-speciation:population}{poblaciones} que se mueven en el mismo sentido, tras el mismo
cambio ambiental y con un mecanismo verosímil: selección. Una única
\omterm{def:g12:selection-drift-speciation:population}{población} diminuta de treinta, sin ninguna causa señalada: la deriva es
la explicación por defecto mientras no se muestre un mecanismo.
\end{solution}

\begin{solution}{exo:g12:selection-drift-speciation:12}
$AA$ $0.85^2 \approx 72\%$, $Aa$ $2 \times 0.85 \times 0.15 \approx
26\%$, $aa$ $0.15^2 \approx 2\%$. Los $aa$ mueren y retiran copias de
$a$; pero los $Aa$, una cuarta parte de la \omterm{def:g12:selection-drift-speciation:population}{población}, sobreviven al
paludismo mejor que los $AA$ y transmiten sus copias de $a$. La pérdida
por los $aa$ se compensa con la ventaja de los $Aa$, y $a$ se mantiene
en el 15\%.
\end{solution}

\begin{solution}{exo:g12:selection-drift-speciation:13}
La elección de pareja por el color: una barrera de comportamiento. En
aguas turbias la barrera falla y las \omterm{def:g10:biodiversity-scales:species}{especies} se hibridan, así que
todavía no se ha acumulado ninguna incompatibilidad genética: la
separación es reciente y aún reversible.
\end{solution}

\begin{solution}{exo:g12:selection-drift-speciation:14}
Alrededor del anillo, cada \omterm{def:g12:selection-drift-speciation:population}{población} se cruza con sus vecinas, con
diferencias pequeñas entre las contiguas; las diferencias se suman a lo
largo del círculo hasta que los extremos, al encontrarse, ya no se
cruzan. Todo el gradiente se ve de una vez: la \omterm{prop:g12:selection-drift-speciation:speciation}{especiación} es cuestión
de pasos pequeños acumulados. A los dos extremos se los llama \omterm{def:g10:biodiversity-scales:species}{especies}
porque, donde se encuentran, se comportan como \omterm{def:g10:biodiversity-scales:species}{especies}: no hay
intercambio de \omterm{def:g10:universal-dna:gene}{genes}.
\end{solution}

\begin{solution}{exo:g12:selection-drift-speciation:15}
La selección hace que los \omterm{def:g11:genes-and-disease:genetic}{portadores} de ciertos \omterm{def:g10:universal-dna:gene}{alelos} se reproduzcan
más \emph{aquí y ahora}: mejores dejando descendientes en este
ambiente, comparados con los demás miembros de la \omterm{def:g12:selection-drift-speciation:population}{población} --- y no
mejores en ningún sentido absoluto. Cuando el ambiente cambia, la
dirección se invierte (las mariposas); la deriva hace que las
\omterm{def:g12:selection-drift-speciation:population}{poblaciones} difieran sin razón de adaptación; la selección sexual
favorece adornos que estorban a la supervivencia. La evolución no tiene
una dirección de mejora, solo una clasificación local.
\end{solution}

\begin{solution}{pb:g12:selection-drift-speciation:1}
\textbf{1.} Copias de $D$: $2 \times 720 + 960 = 2400$; de $d$: $960 +
2 \times 320 = 1600$; de 4000: $p = 0.6$, $q = 0.4$.

\textbf{2.} Esperados $0.36 \times 2000 = 720$, $0.48 \times 2000 =
960$, $0.16 \times 2000 = 320$: exactamente el censo. Sí.

\textbf{3.} $960/1680 \approx 57\%$.

\textbf{4.} Casi todas las copias de $d$ están en los \omterm{def:g11:genes-and-disease:genetic}{portadores}
oscuros. Tras la retirada: 1680 ratones, 2400 copias de $D$ y 960 de
$d$: $q = 960/3360 \approx
0.29$ --- una caída desde 0,40, pero muy
lejos de la eliminación.

\textbf{5.} El gran tamaño: una \omterm{def:g12:selection-drift-speciation:population}{población} insular pequeña deriva, así
que las \omterm{def:g12:selection-drift-speciation:population}{frecuencias} vagan por azar (y no hay migración desde el
continente, y la \omterm{def:g11:mutations:mutation}{mutación} es despreciable).

\textbf{6.} $DD$ 360, $Dd$ 480, $dd$ 320: 1160 reproductores.

\textbf{7.} $D$: $720 + 480 = 1200$; $d$: $480 + 640 = 1120$; de 2320:
$p \approx 0.52$, $q \approx 0.48$.

\textbf{8.} Generación siguiente (por 2000): $DD$ $0.52^2 \approx 27\%$
(535), $Dd$ $\approx 50\%$ (999), $dd$ $\approx 23\%$ (466). Tras los
búhos: 267, 500, 466. $D$: $534 + 500 = 1034$; $d$: $500 + 932 =
1432$;
$q \approx 0.58$.

\textbf{9.} Con $q = 0.58$: $DD$ 18\% (353), $Dd$ 49\% (974), $dd$ 34\%
(673); supervivientes 176, 487, 673; $d$: $487 + 1346 = 1833$ de 2672:
$q \approx 0.69$. Tendencia: 0,40, 0,48, 0,58, 0,69 --- una subida de
alrededor de 0,1 por generación.

\textbf{10.} Todo \omterm{def:g11:genes-and-disease:genetic}{portador} de $D$ es oscuro y queda expuesto a los
búhos, así que la selección actúa sobre todas las copias de $D$ en cada
generación; $q$ sube de forma sostenida. Un \omterm{def:g10:universal-dna:gene}{alelo} \omterm{def:g11:genes-and-disease:genetic}{recesivo} se
escondería en los heterocigotos en cuanto fuera raro, pero uno
\omterm{def:g11:genes-and-disease:genetic}{dominante} no se esconde nunca: $D$ puede llevarse hasta cero.

\textbf{11.} $D$: $4 + 3 = 7$; $d$: $3 + 2 = 5$; $q = 5/12 \approx
0.42$, cercano al 0,40 de la isla por suerte del reparto.

\textbf{12.} Entre los reproductores, $D$ tiene 5 copias y $d$ 1: $q$
en la generación siguiente es $1/6$ de media, pero la única copia de
$d$ puede no pasar a ninguna cría ($q = 0$) o pasar a varias; el
\omterm{def:g11:enzymes-and-phenotype:phenotype}{genotipo} $dd$ ya ha desaparecido, y con él buena parte de la variación.

\textbf{13.} Con tan pocos reproductores, qué \omterm{def:g10:universal-dna:gene}{alelos} se transmiten es
cuestión de azar en cada generación; un \omterm{def:g10:universal-dna:gene}{alelo} raro puede perderse en
una sola camada desafortunada.

\textbf{14.} $d$ derivará y, en algunas generaciones, quedará fijado o
perdido --- cualquiera de las dos cosas por azar. Un islote vecino, con
el mismo punto de partida, puede acabar con el \omterm{def:g10:universal-dna:gene}{alelo} contrario: las dos
islas difieren en el color del pelaje sin ninguna razón adaptativa.

\textbf{15.} La \omterm{prop:g12:selection-drift-speciation:drift}{deriva genética}; el tamaño minúsculo de la colonia.

\textbf{16.} Sí: rara vez se cruzan con los ratones de la isla y los
híbridos son estériles --- aislamiento reproductivo.

\textbf{17.} El mar, una barrera geográfica; la selección en un
ambiente distinto (sin búhos y con otras condiciones) y la \omterm{prop:g12:selection-drift-speciation:drift}{deriva
genética} en una \omterm{def:g12:selection-drift-speciation:population}{población} pequeña.

\textbf{18.} Los ratones que crían en épocas distintas no llegan a
encontrarse como pareja: aun sin ninguna barrera geográfica, no fluyen
\omterm{def:g10:universal-dna:gene}{genes} y las \omterm{def:g12:selection-drift-speciation:population}{poblaciones} siguen divergiendo.

\textbf{19.} Ahora no: están aislados por el comportamiento y por la
esterilidad, y seguirían siendo dos \omterm{def:g10:biodiversity-scales:species}{especies} una junto a la otra. Al
cabo de solo 100 años se habrían cruzado sin trabas y habrían vuelto a
fundirse en una sola \omterm{def:g12:selection-drift-speciation:population}{población}.

\textbf{20.} $q$ pasa de 0,40 a alrededor de 0,69 tras tres
generaciones de búhos; en el islote, $d$ quedó fijado o perdido por
\omterm{prop:g12:selection-drift-speciation:drift}{deriva genética}; dos \omterm{def:g10:biodiversity-scales:species}{especies} al final.
\end{solution}
